% Options for packages loaded elsewhere
\PassOptionsToPackage{unicode}{hyperref}
\PassOptionsToPackage{hyphens}{url}
%
\documentclass[11pt]{article}
\usepackage[margin=1in,headheight=22pt,headsep=20pt,footskip=30pt]{geometry}
\usepackage{amsmath,amssymb}
\usepackage{iftex}
\ifPDFTeX
  \usepackage[T1]{fontenc}
  \usepackage[utf8]{inputenc}
  \usepackage{textcomp} % provide euro and other symbols
\else % if luatex or xetex
  \usepackage{unicode-math} % this also loads fontspec
  \defaultfontfeatures{Scale=MatchLowercase}
  \defaultfontfeatures[\rmfamily]{Ligatures=TeX,Scale=1}
\fi
\usepackage{lmodern}
\ifPDFTeX\else
  % xetex/luatex font selection
\fi
% Use upquote if available, for straight quotes in verbatim environments
\IfFileExists{upquote.sty}{\usepackage{upquote}}{}
\IfFileExists{microtype.sty}{% use microtype if available
  \usepackage[]{microtype}
  \UseMicrotypeSet[protrusion]{basicmath} % disable protrusion for tt fonts
}{}
\makeatletter
\@ifundefined{KOMAClassName}{% if non-KOMA class
  \IfFileExists{parskip.sty}{%
    \usepackage{parskip}
  }{% else
    \setlength{\parindent}{0pt}
    \setlength{\parskip}{6pt plus 2pt minus 1pt}}
}{% if KOMA class
  \KOMAoptions{parskip=half}}
\makeatother
\usepackage{xcolor}
\usepackage{graphicx}
\usepackage{float}
\usepackage{fancyhdr}
\definecolor{GKBlue}{HTML}{18324B}
\definecolor{GKAccent}{HTML}{D47736}
\definecolor{GKInk}{HTML}{263238}
\definecolor{GKMuted}{HTML}{61717D}
\definecolor{GKLight}{HTML}{F1F4F6}
\makeatletter
\def\maxwidth{\ifdim\Gin@nat@width>\linewidth\linewidth\else\Gin@nat@width\fi}
\def\maxheight{\ifdim\Gin@nat@height>\textheight\textheight\else\Gin@nat@height\fi}
\makeatother
% Scale images if necessary, so that they will not overflow the page
% margins by default, and it is still possible to overwrite the defaults
% using explicit options in \includegraphics[width, height, ...]{}
\setkeys{Gin}{width=\maxwidth,height=\maxheight,keepaspectratio}
% Set default figure placement to htbp
\makeatletter
\def\fps@figure{htbp}
\makeatother
\setlength{\emergencystretch}{3em} % prevent overfull lines
\providecommand{\tightlist}{%
  \setlength{\itemsep}{0pt}\setlength{\parskip}{0pt}}
\setcounter{secnumdepth}{-\maxdimen} % remove section numbering
\ifLuaTeX
  \usepackage{selnolig}  % disable illegal ligatures
\fi
\usepackage{bookmark}
\IfFileExists{xurl.sty}{\usepackage{xurl}}{} % add URL line breaks if available
\urlstyle{same}
\hypersetup{
  pdftitle={Geonkick User Guide},
  colorlinks=true,
  linkcolor=GKBlue,
  urlcolor=GKBlue,
  citecolor=GKBlue,
  pdfcreator={LaTeX via pandoc}}

\setlength{\parindent}{0pt}
\setlength{\parskip}{0.55em plus 0.1em minus 0.05em}
\setlength{\emergencystretch}{4em}
\setlength{\itemsep}{0.25em}
\setlength{\topsep}{0.4em}
\setcounter{tocdepth}{2}

\pagestyle{fancy}
\fancyhf{}
\fancyhead[L]{\small\sffamily\bfseries\color{GKBlue}GEONKICK}
\fancyhead[R]{\small\sffamily\color{GKMuted}USER GUIDE}
\fancyfoot[C]{\small\sffamily\color{GKMuted}\thepage}
\renewcommand{\headrulewidth}{0.4pt}
\renewcommand{\footrulewidth}{0pt}
\renewcommand{\headrule}{\hbox to\headwidth{\color{GKAccent}\leaders\hrule height \headrulewidth\hfill}}

\makeatletter
\renewcommand{\section}{\@startsection{section}{1}{\z@}%
  {-2.8ex plus -1ex minus -.2ex}{1.3ex plus .2ex}%
  {\normalfont\Large\bfseries\sffamily\color{GKBlue}}}
\renewcommand{\subsection}{\@startsection{subsection}{2}{\z@}%
  {-2.2ex plus -.8ex minus -.2ex}{0.9ex plus .2ex}%
  {\normalfont\large\bfseries\sffamily\color{GKBlue}}}
\renewcommand{\subsubsection}{\@startsection{subsubsection}{3}{\z@}%
  {-1.8ex plus -.6ex minus -.2ex}{0.7ex plus .2ex}%
  {\normalfont\normalsize\bfseries\sffamily\color{GKBlue}}}
\makeatother

\title{Geonkick User Guide}
\author{Iurie Nistor}
\date{}

\begin{document}
\hypersetup{pageanchor=false}
\makeatletter
\begin{titlepage}
\thispagestyle{empty}
\color{GKInk}
\vspace*{1.2cm}
\noindent\textcolor{GKAccent}{\rule{3.2cm}{3pt}}

\vspace{1.4cm}
\begin{flushleft}
{\sffamily\bfseries\fontsize{32}{38}\selectfont\color{GKBlue}\@title\par}
\vspace{0.5cm}
{\sffamily\Large\color{GKMuted}Author: Iurie Nistor\par}
\end{flushleft}

\vspace{1.3cm}
\noindent\colorbox{GKLight}{%
  \parbox{\dimexpr\textwidth-2\fboxsep\relax}{%
    \vspace{0.35cm}
    \textbf{License:} This work (including all images it uses) is released
    under
    \href{https://creativecommons.org/publicdomain/zero/1.0/deed.en}{CC0 1.0
    Universal (CC0 1.0) Public Domain Dedication}.\par
    \vspace{0.25cm}
    \textbf{Version:} 3.9.0\par
    \vspace{0.25cm}
    \textbf{Website:} \url{https://quamplex.com/geonkick}
    \vspace{0.35cm}
  }}
\vfill
\noindent\textcolor{GKAccent}{\rule{\textwidth}{1pt}}
\end{titlepage}
\makeatother
\hypersetup{pageanchor=true}
\pagenumbering{roman}
\small
\tableofcontents
\normalsize
\clearpage
\pagenumbering{arabic}

\section{Introduction}\label{introduction}

Geonkick is a \href{https://www.gnu.org/philosophy/free-sw.en.html}{free
software} percussive synthesizer. It can synthesize various percussive
sounds, including kicks, snares, hi-hats, shakers, claps, and other
unique sound effects. The synthesizer can also play pre-existing samples
and offers features such as pitch modification, allowing users to modify
and blend samples and synthesis. Geonkick is released under the GNU
General Public License version 3 (and later), developed in C and C++,
and utilizes the CMake build system. It can be used as a standalone
application and as a LV2 \& VST3 plugin. It can run on GNU/Linux and
Windows.

\section{The Sound Designer Craft}\label{the-sound-designer-craft}

Geonkick will not magically improve your percussion sounds. Like with
any other synthesizer, the most important thing is the sound designer's
craft and years of experience. Geonkick only provides the tools
necessary to design the sound. Even I, as a developer, could not design
a good kick. For example, \href{https://haraldrevery.bio.link/}{Harald
Revery}, with his experience in designing percussive sounds,
\href{https://www.youtube.com/watch?v=WnxUgDoOriQ}{shows} how you can
actually design a high-quality kick with Geonkick. You can try Harald
Revery's kick presets that were included in the Geonkick project.

\section{How it works}\label{how-it-works}

Geonkick synthesizes the full instrument length sound into a buffer in
the memory whenever the user modifies the controls or the envelopes. The
resulting instrument sound is then transferred to a module responsible
for playing the instrument. There is no synthesis occurring when the
user presses a key; in other words, if no controls are altered, Geonkick
functions like a sampler. However, this behavior results in the same
real-time synthesis experience for the user as common synthesizers.

Additionally, there is a list of synthesized instruments that enables
Geonkick to function as an instrument kit with specific key mappings.
The Instrument UI (main window) essentially acts as an ``editor'' for
the selected instrument from the kit list.

Each instrument in the kit list is monophonic. However, the instrument
can be directed to the audio output of any of the 16 stereo audio output
channels. Given that the instrument is monophonic, the left and right
channels of a stereo channel are identical.

\clearpage
\section{Block diagram}\label{block-diagram}

\begin{figure}[H]
\centering
\includegraphics[width=0.8\linewidth]{diagram.png}
\caption{Diagram}
\end{figure}

\section{Instrument}\label{instrument}

An instrument is a percussive sound generated by the Geonkick
synthesizer. Its maximum length can be 4 seconds. This sound is played
by Geonkick when an associated MIDI key is pressed.

Geonkick can handle up to 16 instruments in parallel, making it a
multi-timbral synthesizer. There is always a currently selected
instrument that can be modified using the instrument UI. The user can
select any instrument from the kit list and modify it.

\subsection{Playing an instrument}\label{playing-an-instrument}

\begin{itemize}
\tightlist
\item
  Press the \textbf{Play} button in the top bar.
\item
  Through MIDI events. This supports the following playback modes:
  \textbf{Full Length}, \textbf{Note Off}, and \textbf{Cut}.
\end{itemize}

\subsection{Instrument Layers}\label{instrument-layers}

An instrument consists of three layers that can be turned on/off using
the ``E'' buttons in the layer controls. Each layer contains three
oscillators. The sum of these layers is passed to the global controls.
The contribution of each layer to the output can be controlled using the
layer limiter. Although only one layer is controllable through the UI at
a time, the user can switch between layers by clicking an item in the
layer list in the Synth UI.

\subsection{Synthesizer UI}\label{synthesizer-ui}

The Synthesizer's user interface serves as an editor for the current
instrument being synthesized.

\subsection{Top Bar}\label{top-bar}

\begin{itemize}
\tightlist
\item
  \textbf{Settings:} Opens the settings dialog.
\item
  \textbf{Play button:} Plays the current instrument. This does not
  support the playback modes; it always plays the instrument in
  \textbf{Full Length} mode.
\item
  \textbf{Synth / Kit:} Switches between the instrument synthesizer
  controls and the Kit UI.
\item
  \textbf{Reset button:} Resets the current instrument to its default
  control values.
\item
  \textbf{Preset navigator:} Allows selecting factory presets.
\item
  \textbf{Instrument name:} Displays the editable instrument name.
\item
  \textbf{Key:} Displays the current MIDI input key for the instrument.
  Clicking the value label opens a popup for selecting a MIDI key for
  the instrument.
\item
  \textbf{Out:} Selects the audio output channels for the current
  instrument.
\item
  \textbf{MIDI channel:} Selects the MIDI channel for the current
  instrument.
\item
  \textbf{Playback mode button:} Cycles through \textbf{Full Length},
  \textbf{Note Off}, and \textbf{Cut} playback modes.
\item
  \textbf{Tune checkbox:} Toggles tuning for the current instrument.
  When enabled, the instrument's output sound is tuned according to the
  pressed MIDI key.
\end{itemize}

\subsection{Shortcut Keys}\label{shortcut-keys}

\begin{itemize}
\tightlist
\item
  \textbf{`Ctrl + k':} Plays the instrument sound.
\item
  \textbf{`Ctrl + r':} Resets the current instrument to default.
\item
  \textbf{`Ctrl + h':} Hides the envelope, showing only the graph.
\item
  \textbf{`Ctrl + c':} Copies the currently selected instrument.
\item
  \textbf{`Ctrl + v':} Pastes the copied instrument into the currently
  selected instrument.
\item
  \textbf{`Left Ctrl + f':} Scales up the UI by a factor of x1.25, x1.5,
  x1.75, 2x, or restores the original size.
\item
  \textbf{`Left Ctrl + left mouse double-click on an envelope point':}
  Opens the envelope point context input for its value.
\end{itemize}

\textbf{Note:} When running as a plugin, ensure that the focus is on the
main Geonkick window and that the key inputs are enabled by the host for
the plugin. Not all DAWs support this feature reliably.

\subsection{Instrument Envelope Area}\label{instrument-envelope-area}

The instrument envelope area displays:

\begin{itemize}
\tightlist
\item
  \textbf{Instrument sound waveform:} Shows the real-time synthesized
  instrument sound wave form.
\item
  \textbf{Envelopes:} Various types of envelopes displayed on top of the
  instrument sound graph.
\item
  \textbf{Envelope points:} Can be added, removed, and moved by the
  user. Display the values of the envelopes points and also have a
  context input for the value (Ctrl + left mouse double-click on the
  point).
\item
  \textbf{``Bz'' button:} Enables Bezier mode for drawing envelopes.
\item
  \textbf{Envelope Path Label:} Shows the current selected envelope type
  and category. For example, ``L1/OSC1/AMPL'' means the amplitude
  envelope of Layer 1 of Oscillator 1.
\item
  When hovering over/selecting an envelope point, it displays the
  point's coordinates (value, time) in the bottom right corner area.
\item
  Scroll up/down to zoom in/out the envelope area along the time axis.
  Also, drag left/right to move the zoomed area.
\item
  A label ``1:x'' that shows the zoom factor.
\end{itemize}

\subsection{Envelope Points}\label{envelope-points}

Users can add, remove, or move an envelope point. The default first and
last points cannot be removed, only moved.

\begin{itemize}
\tightlist
\item
  \textbf{Left click:} Select and move the point.
\item
  \textbf{Left double click:} Add a new point.
\item
  \textbf{Right click:} Remove the point.
\end{itemize}

\subsection{Bézier Curve for
Envelopes}\label{buxe9zier-curve-for-envelopes}

To add a Bézier curve, you need to insert a control point between two
node points using Left Ctrl + double-click, or press ``Bz'' button in
the bottom of the waveform.

Note: If one of the adjacent points is already a control point, a normal
node point will be added instead. You can only add a control point
between two existing node points.

\subsection{Envelope Point Label}\label{envelope-point-label}

Envelope points have labels showing the current value related to
vertical axes (amplitude, frequency, etc.). For frequency envelope or
cutoff, an additional point label maps the current frequency to a MIDI
note.

\emph{Note: Point key labels are shown for a frequency region close to
or matching the particular key, not just an exact frequency. The exact
value is displayed in the bottom-right corner.}

\subsection{Manually Input Envelope Point
Value}\label{manually-input-envelope-point-value}

Users can set the value of an envelope point manually with a maximum
precision of 4 digits after the floating point. For oscillator frequency
and filter cutoff envelopes, the value can also be entered as note name
from A0 to G\#9 (for example, \texttt{C\#4}, it supports only `\#', not
`b').

\begin{enumerate}
\def\labelenumi{\arabic{enumi}.}
\item
  Press \texttt{Left\ Ctrl} and double-click on the envelope point. An
  input context line edit will be displayed.
\item
  Enter the desired value and press \texttt{Enter}.
\item
  To cancel, press \texttt{Esc} or click outside the line edit widget to
  close it.
\end{enumerate}

\subsection{Display Current Selected Coordinates of Envelope
Point}\label{display-current-selected-coordinates-of-envelope-point}

A more exact values of the coordinates of the currently hovered or
selected envelope point are shown in the bottom-right area, displaying
both value and time in milliseconds.

\subsection{Instrument Wave Form}\label{instrument-wave-form}

The instrument waveform shows the real-time synthesized instrument
sound. Whenever the user changes instrument controls, Geonkick generates
a new graph and updates it.

If the sound exceeds the range of -1.0 to 1.0, the sound is hard-limited
to these values, visible in the instrument graph.

\emph{Issues:} For strictly periodic signals with high frequencies that
cause the graph to be dense, aliasing might occur for specific
frequencies. However, this issue does not affect the played sound or the
exported sound itself.

\subsection{Instrument Preset}\label{instrument-preset}

The instrument preset is in JSON format and uses the file extension
``*.gkit''.

\subsection{Oscillator}\label{oscillator}

Each layer consists of three oscillators, the outputs of which are mixed
with a noise generator to produce the layer's output. Controls related
to oscillators are grouped and arranged vertically.

\begin{itemize}
\tightlist
\item
  \textbf{``OSC1'', ``OSC2'', and ``OSC3'':} Checkboxes to toggle the
  oscillators on/off.
\item
  \textbf{``OSC1-\textgreater OSC2'':} Checkbox that allows Oscillator 1
  to FM modulate Oscillator 2.
\item
  \textbf{Waveform Types:} ``Sine,'' ``Square,'' ``Triangle,''
  ``Sawtooth,'' ``Noise,'' ``Sample'' -- options for wave functions of
  the oscillator.
\item
  \textbf{Noise:} The oscillator will generate noise (``w'' -- white,
  ``b'' -- brownian). Also, there is a ``seed'' spinbox to set the noise
  seed.
\item
  \textbf{``Sample'':} Represents a sampler with a maximum length of 4
  seconds.
\item
  \textbf{``Phase'':} Adjusts the initial phase of the oscillator. For
  samples, this moves the start of the sample up to 1/2 of the
  instrument's sound length.
\item
  \textbf{Amplitude Knob:} Sets the maximum current amplitude of the
  oscillator.
\item
  \textbf{Frequency Knob:} Sets the maximum current frequency of the
  oscillator.
\item
  \textbf{Pitch Knob:} Sets the pitch shift range when the waveform is
  set as a sample.
\item
  \textbf{Density Knob:} When the noise wave function is selected, this
  knob controls the density of the noise.
\item
  \textbf{Pitch Button:} Shows the envelopes for the pitch shift.
\item
  \textbf{Density Button:} Shows the envelope of the noise density.
\item
  \textbf{``Filter''} and \textbf{``Distortion''} effects buttons:
  Buttons to enable and switch to the filter or distortion effect UI.
\item
  \textbf{Distortion:}
\item
  \textbf{Filter:} a digital state-variable filter

  \begin{itemize}
  \tightlist
  \item
    Cutoff knob and \textbf{Cutoff} envelope button, filter types:
    \textbf{LP} (low-pass), \textbf{BP} (band-pass), \textbf{HP}
    (high-pass), \textbf{Resonance} knob, and \textbf{Q} resonance
    envelope button.
  \end{itemize}
\end{itemize}

\textbf{Note:} Loading/saving presets preserves the noise pattern, seed,
and density pattern for the oscillator.

\subsection{Global Controls}\label{global-controls}

Global controls are applied after the layers mixer. Arranged vertically
and applied from top to bottom.

\begin{itemize}
\tightlist
\item
  \textbf{Amplitude Knob:} Controls the maximum amplitude after the
  layers mixer.
\item
  \textbf{``Amplitude'' Button:} Shows the general amplitude envelope.
\item
  \textbf{``Length'':} Sets the maximum sound length of the instrument.
\item
  \textbf{Filter:} Similar to oscillator.
\item
  \textbf{Distortion:} Similar to oscillator.
\item
  \textbf{Humanizer:} Velocity and timing humanizer.
\end{itemize}

\subsection{Limiter}\label{limiter}

The limiter control is applied to the audio output of the instrument,
controlling only the currently selected instrument.

\subsection{Filter}\label{filter}

Geonkick implements a digital state-variable filter applied to
oscillators, and after the layers mixer (instrument global filter).

Parameters:

\begin{itemize}
\tightlist
\item
  \textbf{``Filter'' Buttons:} Enable/Disable filter.
\item
  \textbf{Cutoff Knob:} Controls the cutoff frequency of the filter in a
  logarithmic range (20Hz - 20kHz).
\item
  \textbf{``Cutoff'' Button:} Shows the cutoff envelope.
\item
  \textbf{Resonance Knob:} Logarithmic range from 0.01 to 10.
\item
  \textbf{``Q'' Button:} Shows the envelope for the resonance.
\item
  \textbf{Filter Type Buttons:} ``LP'' for low-pass, ``BP'' for
  band-pass, and ``HP'' for high-pass filters.
\end{itemize}

The filter becomes unstable for extreme combinations of cutoff frequency
and Q parameter. If it becomes unstable, the instrument sound will be
hard-limited. Proper adjustments of the Q factor are necessary to
maintain desired cutoff frequency ranges. But this behavior can be used
as a feature to generate special effect sounds.

\subsection{Distortion}\label{distortion}

The distortion effect (for oscillators and global) is applied after the
filter and has the following controls:

\begin{itemize}
\tightlist
\item
  \textbf{Distortion Type:} Functions used to create distortion:
  \emph{hard clip}, \emph{tangent}, \emph{arctangent},
  \emph{exponential}, \emph{logarithm}, and \emph{polynomial}.
\item
  \textbf{Drive} knob,
\item
  \textbf{Drive envelope Button},
\item
  \textbf{Input} and \textbf{Output} Limiters.
\end{itemize}

\subsection{Tune}\label{tune}

In the top bar, the ``Tune'' checkbox toggles tuning for the currently
selected instrument. When enabled, Geonkick will stop mapping the
instrument to a specific key and will tune the sound relative to the
untuned key A4 (in number of semitones).

\section{Kit UI}\label{kit-ui}

Geonkick is multi-instrument by default. Each synthesized instrument is
added to a list, forming a kit of up to 16 instruments. To access the
kit UI, click the \textbf{``Kit''} button on the top bar.

\subsection{Default Kit State}\label{default-kit-state}

When the user opens a new Geonkick instance, it will have a single
default instrument in the kit UI list. The key mapping for this
instrument will be set to any key, and the channel of the default
instrument will be directed to channel number 0.

\section{Kit UI}\label{kit-ui-1}

The instrument list has a maximum limit of 16 instruments and displays
the instrument names and their MIDI key mappings in a grid. The Kit UI
shows the currently selected instrument from the list, which can be
controlled. In the Kit list, the current instrument is highlighted.

\begin{itemize}
\tightlist
\item
  \textbf{``+'' button (plus icon):} Adds a new instrument.
\item
  \textbf{``-'' button (minus icon):} Removes the current instrument.
\item
  \textbf{Up/down arrow buttons:} Move the currently selected instrument
  up or down.
\item
  \textbf{Copy:} Copy instrument to clipboard.
\item
  \textbf{Paste:} Paste instrument from clipboard.
\item
  \textbf{Duplicate:} Duplicate an instrument (won't duplicate if there
  are already maximum number of instruments created)
\item
  \textbf{Left click:} Selects an instrument.
\item
  \textbf{Left double-click on the instrument name:} Edits the
  instrument name.
\item
  \textbf{Left double-click on the instrument waveform:} Plays the
  instrument.
\item
  \textbf{MIDI channel spinbox:} Sets the MIDI input channel for the
  instrument.
\item
  \textbf{Note:} Selects the MIDI note name.
\item
  \textbf{Oct:} Selects the MIDI note octave.
\item
  \textbf{Playback mode button:} Cycles through \textbf{Full Length},
  \textbf{Note Off}, and \textbf{Cut} playback modes.
\item
  \textbf{Choke:} Selects the choke group for the instrument.
\item
  \textbf{Limiter slider:} Controls the limiter for the current
  instrument.
\item
  \textbf{``M'' button:} Mutes the current instrument.
\item
  \textbf{``S'' button:} Solos the current instrument.
\item
  \textbf{Output channel:} Selects the audio output channel for the
  instrument.
\item
  \textbf{Reset:} Resets the kit, and adds only one default instrument
\end{itemize}

\subsection{Current Instrument}\label{current-instrument}

When the user selects an instrument from the Kit list, it becomes the
current instrument controlled by both the Kit and Synthesizer UIs.

\subsection{Instrument vs Channel}\label{instrument-vs-channel}

In Geonkick's Kit UI, instruments and channels are distinct entities. An
instrument from the kit UI can be redirected to any audio output channel
(from channels 1 to 16). This allows multiple instruments to be directed
to a single channel. However, assigning multiple channels to a single
instrument is not possible from the Kit UI (this feature may be provided
by the host).

\section{Geonkick Preset}\label{geonkick-preset}

Geonkick supports two JSON preset formats:

\begin{itemize}
\tightlist
\item
  \textbf{Instrument presets} use the ``*.gkick'' extension and contain
  one instrument, along with preset metadata such as its name, author,
  and URL. Loading one replaces only the currently selected instrument.
  When saved inside a kit, the instrument's preset metadata is not
  retained; the kit's metadata applies to the complete preset.
\item
  \textbf{Kit presets} use the ``*.gkit'' extension and contain one or
  more instruments, along with kit-level metadata such as its name,
  author, and URL.
\end{itemize}

Both formats can be opened and saved from the \textbf{Preset Browser}.
Enter a filename ending in ``\emph{.gkick'' to save the selected
instrument, or ''}.gkit'' to save the whole kit.

\section{Sidebar}\label{sidebar}

The sidebar contains the \textbf{Preset Browser} and \textbf{Sample
Browser}. It can be shown or hidden from the settings, and this applies
to all Geonkick instances. Use the top menu of the sidebar to switch
between the \textbf{Preset Browser} and \textbf{Sample Browser}.

Both the Preset and Sample browsers have a top menu with the following
controls:

\begin{itemize}
\tightlist
\item
  \textbf{``Forward'' \& ``Backwards''} --- Navigate through browsing
  history.
\item
  \textbf{``+''} --- Create a new directory in the current path.
\item
  **``*``** --- Bookmark/unbookmark the current path.
\item
  \textbf{Bookmarks bar} --- Contains bookmarked paths, as well as
  default paths like factory presets and user paths.
\item
  \textbf{Breadcrumb bar} --- Displays the current path as clickable
  segments for easy navigation.
\item
  \textbf{Save/Export button} --- In the Preset Browser, saves the
  current preset. In the Sample Browser, exports the current
  instrument's waveform.
\end{itemize}

\subsection{Preset Browser}\label{preset-browser}

The Preset Browser includes a list that shows the contents of the
currently selected preset folder. When the user selects a path that
contains presets, they will be displayed in the list.

The lock button next to Save toggles \textbf{Keep instrument settings}
when loading an instrument preset from either the Preset Browser or the
preset navigator. When enabled, the preset does not change the
instrument's MIDI key, MIDI channel, output channel, choke group, or
play mode. The setting is off by default and its state is remembered
between sessions. Kit presets still load as complete kits.

\subsection{Sample Browser}\label{sample-browser}

The Sample Browser has additional controls in the bottom bar:

\begin{itemize}
\tightlist
\item
  \textbf{``O1, O2, O3''} - When a sample file is selected, pressing one
  of these buttons will load the sample into the corresponding
  oscillator (O1, O2, or O3), and the oscillator's waveform will be set
  to ``sample''.
\item
  \textbf{``Export as''} - A spinbox where the export format can be
  selected: WAV, FLAC, OGG, or SFZ.
\item
  \textbf{``M''}, \textbf{``S''} - Export as stereo or mono (not
  applicable for SFZ format).
\end{itemize}

\subsection{Export}\label{export}

The export feature can be accessed by clicking the \textbf{``Save''}
button in the \textbf{Sample Browser}.

\begin{itemize}
\tightlist
\item
  \textbf{File format:} ``FLAC16'', ``FLAC24'', ``WAV16'', ``WAV24'',
  ``WAV32'', ``Ogg'', ``SFZ''
\end{itemize}

Geonkick will export only the currently selected instrument.

\subsection{Settings}\label{settings}

Settings can be accessed from the top bar. The settings widget allows
the user to configure various settings. These settings apply to all
Geonkick instances. For settings to take effect (especially UI-related
ones), it may be necessary to restart the UI or even the instance
(restart the host, reload the plugin, or reopen the standalone
application).

\begin{itemize}
\tightlist
\item
  \textbf{Force all MIDI channels to} - Forces all instrument MIDI input
  channels to a specific channel (applies to all instances).
\item
  \textbf{Scale UI} - Sets the scale of the UI (options: no scale,
  1.25x, 1.5x, 1.75x, and 2x). Useful for HiDPI screens.
\item
  \textbf{Show sidebar} - Toggles the visibility of the sidebar.
\end{itemize}

\section{LV2 \& VST Plugin}\label{lv2-vst-plugin}

There are two versions of the plugin: the default version and the
``Single'' version.

\begin{itemize}
\item
  The default plugin version supports kits with 16 instruments and
  creates 16 stereo audio output channels per plugin instance.
\item
  The ``Single'' plugin version does not support kits and only includes
  one instrument and one stereo audio output channel.
\item
  VST3 plugin has only the default version (kit version).
\end{itemize}

\section{Standalone}\label{standalone}

The standalone version of Geonkick offers the same features as the LV2
or VST3 plugin.

\section{Build \& Install}\label{build-install}

For instructions on building and installing Geonkick, please refer to
the README file.

\end{document}
